\scanpage{120}

$G_0$ est fini i.e. $S$ fini. Sinon ce n'est pas tellement clair sur les définitions. On peut quand = dire que \deleted{pour} $G_0^+$ est simplement transitif sur l'ens. $\vec A_\lambda$ des arêtes orientées de \deleted{$A$} $P_\lambda$, donc \deleted{\unsure{si on a|si on savait}} pour montrer que $G_0$ est transitif (donc \deleted{deux repères $\unsure{r,r'}$ de $P_\lambda$,} simpl. transitif) sur l'ens des repères de $A_{\unsure{\lambda|0}}$, \deleted{$r_0=(s_0,t_0,u_0)\qquad r_1=(s_1,t_1,u_1)$} \deleted{il} il suffit de prouver que pour \deleted{\unsure{tt}} un repère $(s_0,t_0,u_0)$\inserted{(donc pour \unsure{tt|$\forall$} repère)}, il existe un $u\in G_0$ \deleted{tel} qui fixe $s_0,t_0$, et \uline{non} $u_0$, \deleted{Mais s'il}\inserted{i.e.} qui n'est pas l'identité. Mais s'il n'en était ainsi, cela signifierait que $G_0$ est \uline{simplement} transitif sur l'ens $\vec A_\lambda$ (le stabilisateur d'un $\vec a\inserted{=(s_0,t_0)}\in\vec A_\lambda$ dans $G$ est réduit à l'élément neutre), ce qui contredit le fait que $G_0^+$ est simplement transitif et \deleted{\unsure{d'i. de}} $G_0^+\subsetneq G_0$. Donc on a bien un iso

\[
G_0\mathrel{\overset{\sim}{\rule{2cm}{0.4pt}}}\operatorname{Aut}(P_\lambda)=G_{\lambda\deleted{\unsure{\lambda|?}}},
\tag*{$(43)$}
\]

on aimerait prouver que cet iso. induit

\[
G_0^+\xrightarrow{\ \sim\ }G_\lambda^+\qquad ?
\tag*{$(44)$}
\]
% archive sheet number: 119
